\section*{Capítulo \ref{ch:b2:basins-compared} --- Cuencas sedimentarias comparadas}

\begin{solution}{exo:b2:basins-compared:1}
(a) Intracratónica: subsidencia lenta a impulsos durante cientos de millones de años, sin fallas
importantes; su causa se discute. (b) De rift: adelgazamiento mecánico de la corteza durante el estiramiento,
y después contracción térmica al enfriarse la litosfera. (c) De margen pasivo: la subsidencia térmica
heredada de la etapa de rift, amplificada por la carga de una gruesa cuña sedimentaria. (d) y (e)
De antepaís: la litosfera flexionada hacia abajo por la carga del cinturón de cabalgamientos, los Alpes y el Himalaya.
\end{solution}

\begin{solution}{exo:b2:basins-compared:2}
(a) $Y = 3000\times(3330 - 2300)/(3330 - 1030) = 3000\times1030/2300 = \qty{1343}{m}$.
(b) $1343 + 200 = \qty{1543}{m}$. (c) Sin su carga, el basamento estaría a \qty{1343}{m}:
$1 - 1343/3000 = \qty{55}{\percent}$ del espesor corresponde al espacio que abre la propia carga.
\end{solution}

\begin{solution}{exo:b2:basins-compared:3}
(a) $\beta = 35/25 = 1.4$. (b) $125/1.4 = \qty{89.3}{km}$. (c) $S_i = 4334\times(1 - 1/1.4) =
\qty{1238}{m}$.
\end{solution}

\begin{solution}{exo:b2:basins-compared:4}
(a) $\tau = (1.25\times10^{5})^2/(\pi^2\times8\times10^{-7}) = \qty{1.979e15}{s}$, es decir,
$1.979\times10^{15}/3.156\times10^{13} = \qty{62.7}{Myr}$. (b) $62.7\times0.8 = \qty{50.2}{Myr}$.
(c) $62.7\times(100/125)^2 = \qty{40.1}{Myr}$. (d) $1 - e^{-1} = \qty{63}{\percent}$; $1 -
e^{-3} = \qty{95}{\percent}$.
\end{solution}

\begin{solution}{exo:b2:basins-compared:5}
(a) $0.63\,e^{-1000/1960} = 0.378$; $0.63\,e^{-2000/1960} = 0.227$. (b) $T_s = 1000 -
0.63\times1960\times(0.6004 - 0.3605) = \qty{704}{m}$. (c) $m = 2700\times704 + 1030\times296 =
\qty{2.206e6}{kg.m^{-2}}$; $\bar\rho = \qty{2206}{kg.m^{-3}}$.
\end{solution}

\begin{solution}{exo:b2:basins-compared:6}
(a) $z_2' \leftarrow 704 + 1234.8\,(1 - e^{-z_2'/1960})$: 1077, 1226, 1278, 1296, 1301, 1303,
1304; la capa tiene \qty{1304}{m} de espesor. (b) $(2700\times704 + 1030\times600)/1304 =
\qty{1932}{kg.m^{-3}}$. (c) En \qty{304}{m}, o un \qty{30}{\percent}: el agua que la compactación
ha expulsado.
\end{solution}

\begin{solution}{exo:b2:basins-compared:7}
(a) $1 - e^{-100/62.7} = 0.7971$. Para $\beta = 1.5$: $3207\times0.4135\times0.7971 =
\qty{1057}{m}$; para $\beta = 2$: $3207\times0.6366\times0.7971 = \qty{1627}{m}$. (b) Sí: de 1100 a
\qty{1400}{m} queda entre ambos. (c) $\frac{\beta}{\pi}\sin\frac{\pi}{\beta} = 1250/(3207\times0.7971)
= 0.4890$, de modo que $\beta = 1.6 + 0.1\times(0.4890 - 0.4705)/(0.5204 - 0.4705) = 1.64$, dentro del
50 al \qty{100}{\percent} de estiramiento inferido para la cuenca.
\end{solution}

\begin{solution}{exo:b2:basins-compared:8}
(a) La densidad del relleno sustituye a $\rho_w$ en el denominador. (b) El numerador es
$125\,000\times(1/3)\times74.69 = \num{3.112e6}$; $S_i = 3.112\times10^{6}/(3184.4 - 2300) =
\qty{3519}{m}$. (c) Bajo agua, $3.112\times10^{6}/2154.4 = \qty{1445}{m}$: la cuenca rellena de sedimento
es 2.44 veces más profunda.
\end{solution}

\begin{solution}{exo:b2:basins-compared:9}
(a) $2171 - 100\times3330/2300 = 2171 - 145 = \qty{2026}{m}$. (b) $Y$ varía metro a metro
con $W_d$: de 1971 a \qty{2371}{m}. (c) La de la profundidad del agua, $\pm\qty{200}{m}$, frente a
\qty{145}{m} para el nivel del mar.
\end{solution}

\begin{solution}{exo:b2:basins-compared:10}
(a) $S_i$ tiene el signo del corchete de su numerador, ya que el denominador, $3184.4 - 1030$, es
positivo; es negativa, un levantamiento, cuando lo es el corchete. (b) $(\rho_m - \rho_c)t_c/a =
\alpha T_1\rho_m/2$ da $t_c = a\,\alpha T_1\rho_m/(2(\rho_m - \rho_c)) = 125\times0.04372\times
3330/1060 = \qty{17.17}{km}$. (c) $t_c = 17.17/(1 - 0.04372\times17.17/250) = \qty{17.22}{km}$.
(d) Donde la corteza tiene menos de unos \qty{17}{km} antes del estiramiento: en la litosfera oceánica,
o en una corteza continental ya adelgazada por un rifting anterior.
\end{solution}

\begin{solution}{exo:b2:basins-compared:11}
(a) Antes del estiramiento, la columna contiene la litosfera de masa $M = \rho_c t_c(1 - \alpha T_1
t_c/(2a)) + \rho_m(a - t_c)(1 - \alpha T_1(a + t_c)/(2a))$; después, $\rho_w S_i + M/\beta +
\rho_m(1 - \alpha T_1)(a - S_i - a/\beta)$. Igualándolas, $S_i[\rho_m(1 - \alpha T_1) - \rho_w] =
(1 - 1/\beta)[\rho_m(1 - \alpha T_1)a - M]$, y desarrollando $M$, $\rho_m(1 - \alpha T_1)a - M =
(\rho_m - \rho_c)t_c(1 - \alpha T_1 t_c/(2a)) - \rho_m\alpha T_1 a/2$: la fórmula. (b) El perfil final es $T_1 z/a$; el inicial lo supera
en $T_1 z(\beta - 1)/a$ por encima de $a/\beta$ y en $T_1(1 - z/a)$ por debajo. Las integrales son $T_1
a(\beta - 1)/(2\beta^2)$ y $T_1 a(\beta - 1)^2/(2\beta^2)$, cuya suma es $T_1 a(1 - 1/\beta)/2$.
La contracción, $\alpha$ por esta cantidad, cargada de agua, da la expresión. (c) $3.28\times10^{-5}
\times3330\times1333\times125\,000\times0.5/(2\times2300) = \qty{1978}{m}$, frente a $3207\times
2/\pi = \qty{2042}{m}$: el primer término exagera la subsidencia final en \qty{64}{m}, un
\qty{3}{\percent}.
\end{solution}

\begin{solution}{exo:b2:basins-compared:12}
(a) $3207\times0.6366\times(1 - e^{-150/62.7}) = 2042\times0.9086 = \qty{1855}{m}$. (b) $1855\times
2300/(3330 - 2400) = \qty{4588}{m}$. (c) De 6.9 a \qty{13.8}{km}. (d) El modelo deja fuera varias
cosas que podrían hacer más profunda la cuenca: un manto litosférico más estirado que la corteza, una respuesta regional
(flexural) a la carga del sedimento, un levantamiento transitorio del margen durante el rifting
seguido de hundimiento, o un aporte de sedimento tan grande que la cuenca quedó sobrealimentada y cargada
por encima del nivel del mar en la plataforma.
\end{solution}

\begin{solution}{pb:b2:basins-compared:1}
\textbf{1.} $0.49\,e^{-4420/3704} = 0.149$ en el techo; $0.49\,e^{-6030/3704} = 0.096$ en la base.
\textbf{2.} $T_s = 1610 - 0.49\times3704\times(0.3032 - 0.1963) = \qty{1416}{m}$.
\textbf{3.} Lutita: $1110 - 0.63\times1960\times(e^{-3310/1960} - e^{-4420/1960}) = 1110 - 1234.8
\times(0.1847 - 0.1049) = \qty{1011}{m}$.
\textbf{4.} Arena lutítica: $2680\times1004 + 1030\times696 = \qty{3.408e6}{kg.m^{-2}}$; creta:
$2710\times1409 + 1030\times201 = \qty{4.025e6}{kg.m^{-2}}$; lutita: $2700\times1011 + 1030\times99 =
\qty{2.832e6}{kg.m^{-2}}$; arena: $2650\times1416 + 1030\times194 = \qty{3.952e6}{kg.m^{-2}}$. Total
\qty{1.4217e7}{kg.m^{-2}}, $\bar\rho_s = 14\,216\,910/6030 = \qty{2358}{kg.m^{-3}}$.
\textbf{5.} $z_2' \leftarrow 1416 + 1815\,(1 - e^{-z_2'/3704})$: 1993, 2171, 2221, 2235, 2238, 2239,
2239: \qty{2239}{m}.
\textbf{6.} $(2650\times1416 + 1030\times823)/2239 = \qty{2055}{kg.m^{-3}}$.
\textbf{7.} $Y = 2239\times(3330 - 2055)/2300 + 200 = 1241 + 200 = \qty{1441}{m}$.
\textbf{8.} $Y = 6030\times(3330 - 2358)/2300 + 100 = 2548 + 100 = \qty{2648}{m}$. Las unidades se
miden a sus profundidades actuales: ya están tan compactadas como hoy.
\textbf{9.} $(2171 - 1441)/50 = \qty{14.6}{m}$ por millón de años de 150 a \qty{100}{Ma};
$(2453 - 2171)/40 = \qty{7.05}{m}$ de 100 a \qty{60}{Ma}; $(2648 - 2453)/60 = \qty{3.25}{m}$
desde entonces: el hundimiento se ralentiza.
\textbf{10.} $6030 + 100 = \qty{6130}{m}$, 2.31 veces la subsidencia tectónica.
\textbf{11.} $(6130 - 2648)/6130 = \qty{57}{\percent}$.
\textbf{12.} $2171 - 100\times3330/2300 = \qty{2026}{m}$.
\textbf{13.} El punto del final del rifting, \qty{150}{Ma}: $S_i = \qty{1441}{m}$, si no hubo subsidencia
térmica durante los \qty{15}{Myr} del rifting.
\textbf{14.} $\beta = 1/(1 - 1441/4334) = 1/0.6675 = 1.50$.
\textbf{15.} $35/1.50 = \qty{23.3}{km}$ de corteza y $125/1.50 = \qty{83.3}{km}$ de litosfera.
\textbf{16.} $(1.50/\pi)\sin(\pi/1.50) = 0.4775\times0.8660 = 0.4135$; $3207\times0.4135 =
\qty{1326}{m}$.
\textbf{17.} $S_i = 4334/3 = \qty{1445}{m}$. A \qty{100}{Ma}, $t = \qty{50}{Myr}$: $1326\times(1 -
e^{-50/62.7}) = \qty{729}{m}$, $Y = \qty{2173}{m}$; a \qty{60}{Ma}: $1326\times0.7620 =
\qty{1010}{m}$, $Y = \qty{2455}{m}$; hoy: $1326\times0.9086 = \qty{1205}{m}$, $Y = \qty{2650}{m}$.
\textbf{18.} El modelo es \qty{2}{m} más profundo en cada punto: 2173 frente a 2171, 2455 frente a 2453,
2650 frente a 2648.
\textbf{19.} Paleoprofundidades del agua conocidas solo con un margen de cien metros o más; niveles del mar pasados
distintos del actual; un rifting que duró \qty{15}{Myr}, durante el cual parte del calor
ya se había perdido; flexión en lugar de isostasia local; un manto estirado más o menos que la corteza;
una corteza que no tenía \qty{35}{km} de espesor antes del rifting.
\textbf{20.} $1326\times(e^{-150/62.7} - e^{-200/62.7}) = \qty{67}{m}$.
\textbf{21.} $66.6\times2300/(3330 - 2400) = \qty{165}{m}$.
\textbf{22.} $1 - e^{-150/62.7} = \qty{91}{\percent}$.
\textbf{23.} $\tau\ln2 = \qty{43.5}{Myr}$ después del final del rifting: a \qty{106.5}{Ma}.
\textbf{24.} $\beta = \textbf{1.50}$: la corteza se estiró hasta una vez y media su anchura.
\end{solution}
